\emph{Set-up.} Units $\hbar = c = 1$. On one generator of a horizon, with affine parameter $u \in \R$, the modes form the continuum mode space $\CAR(L^{2}(\R, du)\otimes K)$ of Definition \claimref{obj.hawking.continuum-mode-space}, with fields $\psi(u)$, $\{\psi(u), \psi^{*}(v)\} = \delta(u - v)\otimes\id_{K}$. For $a \in \R$ let $\mathfrak{A}_{a} = \CAR(L^{2}((a, \infty), du)\otimes K)$, the modes of the future half beyond $a$, and $\beta^{(a)}_{t}(\psi(u)) = e^{t/2}\psi(a + e^{t}(u - a))$, the boost about $a$ (rapidity $t$, weight $\tfrac12$, the weight that preserves the anticommutator), which maps $\mathfrak{A}_{a}$ onto itself. Its one-particle generator $h_{a}$ is self-adjoint on $L^{2}((a, \infty))\otimes K$, diagonal on the Mellin modes $(u - a)^{-1/2 + I_{4}w}$ with spectrum $\R$ (the unit of the mode space is $I_{4}$; numbers are read in $\mathrm{span}\{1, I_{4}\}$). The null vacuum $\omega_{0}$ is the gauge-invariant quasi-free state with symbol $\chi_{(-\infty, 0)}(P)\otimes\id_{K}$, $P = I_{4}\,d/du$, two-point function $\omega_{0}(\psi^{*}(v)\psi(u)) = \tfrac{I_{4}}{2\pi}(u - v + I_{4}0)^{-1}\otimes\id_{K}$; it is the unique gauge-invariant quasi-free state of (N$_{+}$) (Premise \claimref{synthesis.n-plus-null-ground-state}; Theorem \claimref{synthesis.positive-null-energy-gives-r-premises}(i), which characterises the states satisfying it).

\emph{Step 1: the KMS states about one point.} By Theorem \claimref{hawking.i-state-continuum} with Fact \claimref{hawking.imported-ideal-fermi-gas-uniqueness}, for each $\beta > 0$ there is exactly one state of $\mathfrak{A}_{a}$ that is KMS at $\beta$ for $\beta^{(a)}$, among all states: the quasi-free I-state $\omega^{(a)}_{\beta}$ with symbol $(1 + e^{\beta h_{a}})^{-1}$. Its two-point function is (the Fourier integral and its residues written with $i$ for the unit $I_{4}$ of $\mathrm{span}\{1, I_{4}\}$) $\int\frac{dw}{2\pi}(1 + e^{\beta w})^{-1}(XY)^{-1/2}(X/Y)^{iw}$ with $X = u - a$, $Y = v - a$. Off the diagonal the Fermi transform is a sum of residues at $w = i\pi(2n + 1)/\beta$: with $r = e^{-\pi s/\beta}$, $s = \ln(X/Y) > 0$, the residue series is $\sum_{n \ge 0}r^{2n + 1} = r/(1 - r^{2}) = 1/(2\sinh(\pi s/\beta))$ (checked, with the exact remainder $r^{2N+1}/(1 - r^{2})$ of the partial sums); so, standard,
\[
\omega^{(a)}_{\beta}(\psi^{*}(v)\psi(u)) = \frac{I_{4}}{2\beta}\,\frac{(XY)^{-1/2}}{\sinh(\pi s/\beta)}\otimes\id_{K}\qquad(u \ne v),
\]
the value for $s < 0$ following by the same computation. Distinct periods give distinct states (distinct symbols), and a KMS state is invariant under its dynamics.

\emph{Step 2: the kernel at fixed separation.} Put $d = u - v = X - Y$ and $E = e^{s} = X/Y$; then $X = dE/(E - 1)$, $Y = d/(E - 1)$, $XY = d^{2}E/(E - 1)^{2}$ and $(XY)^{-1/2} = (E - 1)/(d\sqrt{E}) = 2\sinh(s/2)/d$. So, off the diagonal,
\[
\omega^{(a)}_{\beta}(\psi^{*}(v)\psi(u)) = \frac{I_{4}}{\beta d}\,\frac{\sinh(s/2)}{\sinh(\pi s/\beta)}.
\]
Translating $u$ and $v$ together by $\tau$ keeps $d$ and moves $s$ monotonically through an interval. At $\beta = 2\pi$ the ratio is $1$ and the kernel is $\tfrac{I_{4}}{2\pi d}$, the null vacuum's, for every position (checked exactly on $E = q^{4}$ at four values of $q$ and two of $d$). For $\beta \ne 2\pi$ the ratio $\sinh(s/2)/\sinh(\pi s/\beta)$ is not constant on any interval: it is analytic in $s > 0$, it tends to $\beta/2\pi$ as $s \to 0$ and behaves as $e^{(1/2 - \pi/\beta)s}$ times a constant as $s \to \infty$, so constancy on an interval would make it constant on $(0, \infty)$ and force $\pi/\beta = \tfrac12$. So $\omega^{(a)}_{\beta}$ restricted to any $\mathfrak{A}_{a'}$, $a' > a$, is translation invariant exactly when $\beta = 2\pi$ (the check confirms that the kernel takes distinct values at four positions for $\beta = 4\pi, \tfrac43\pi, \pi, \tfrac23\pi, \tfrac12\pi$).

\emph{Step 3: two points fix the period.} The affine maps $D_{a}(\lambda)\colon x \mapsto a + \lambda(x - a)$ satisfy $D_{a}(\lambda)\circ D_{b}(1/\lambda) = T((\lambda - 1)(b - a))$, a translation (checked). Let $a < b$ be two points such that the restriction of a state $\omega$ to $\mathfrak{A}_{a}$ is KMS at some $\beta_{a} > 0$ for $\beta^{(a)}$ and its restriction to $\mathfrak{A}_{b}$ is KMS at some $\beta_{b} > 0$ for $\beta^{(b)}$. By Step 1, $\omega|_{\mathfrak{A}_{a}} = \omega^{(a)}_{\beta_{a}}$, and $\omega|_{\mathfrak{A}_{b}}$ is invariant under $\beta^{(b)}$. On $\mathfrak{A}_{b}$, the boost $\beta^{(b)}$ with $\lambda = 1/\mu < 1$ followed by the boost $\beta^{(a)}$ with $\mu > 1$ (both map $(b, \infty)$ into itself, and $\omega^{(a)}_{\beta_{a}}$ is invariant under the second) is the translation by $(\mu - 1)(b - a) > 0$; every positive translation is reached. So the two-point function of $\omega^{(a)}_{\beta_{a}}$ on $(b, \infty)^{2}$ is invariant under positive translations, and by Step 2, $\beta_{a} = 2\pi$. Then $\omega|_{\mathfrak{A}_{a}} = \omega^{(a)}_{2\pi}$, whose kernel is the null vacuum's, $\tfrac{I_{4}}{2\pi}(u - v + I_{4}0)^{-1}$: it is $\omega_{0}|_{\mathfrak{A}_{a}}$ (Theorem \claimref{synthesis.positive-null-energy-gives-r-premises}(ii) at $a = 0$, carried to every $a$ by the translation invariance of $\omega_{0}$; the weight-$\tfrac12$ dilation invariance of $(u - v)^{-1}$ about every point is checked). Its restriction to $\mathfrak{A}_{b}$ is $\omega_{0}|_{\mathfrak{A}_{b}} = \omega^{(b)}_{2\pi}$, so by uniqueness of the period, $\beta_{b} = 2\pi$ as well. (Proved.)

\emph{Step 4: a set unbounded below fixes the state.} Let $A \subset \R$ be unbounded below, with at least two points, and suppose that for every $a \in A$ the restriction of $\omega$ to $\mathfrak{A}_{a}$ is KMS at some $\beta_{a} > 0$ for $\beta^{(a)}$. By Step 3 (every $a \in A$ has another point of $A$ below it) every $\beta_{a} = 2\pi$ and $\omega = \omega_{0}$ on every $\mathfrak{A}_{a}$, $a \in A$. These algebras increase as $a$ decreases, and $\bigcup_{a \in A}L^{2}((a, \infty))\otimes K$ is dense in $L^{2}(\R)\otimes K$; since $\|a(f) - a(g)\| = \|f - g\|$ in the CAR algebra, $\bigcup_{a \in A}\mathfrak{A}_{a}$ is norm dense, and two states agreeing on a dense subalgebra agree. So $\omega = \omega_{0}$: gauge-invariant, quasi-free, invariant under the null translations with positive generator, the state of (N$_{+}$). Conversely $\omega_{0}$ is KMS at $2\pi$ about every point (Step 3). No quasi-freeness, gauge invariance, stationarity or temperature is assumed. (Proved.) The mirror statement holds for past halves $(-\infty, a)$, $a$ in a set unbounded above, with the boost future-directed there.

\emph{Step 5: what does not suffice.} (a) One point: the stationary states thermal on each half at period $2\pi$ about the edge with no cross-edge correlation are KMS about $0$ on $\mathfrak{A}_{0}$ and are not $\omega_{0}$ (Theorem \claimref{synthesis.stationary-horizon-state-is-its-own-edge-scaling-limit}). (b) Points bounded below: if $A \subset [b, \infty)$, any state equal to $\omega_{0}$ on $\mathfrak{A}_{b}$ and arbitrary on the modes below $b$ (for instance a product with an even state there) satisfies the hypothesis at every $a \in A$; the check exhibits a finite model in which a symbol with the vacuum's compression on every $[a, \infty)$, $a \ge b$, differs from it. (c) A common period other than $2\pi$ at two points: impossible by Step 3.

\emph{Step 6: de Sitter horizons (argued).} On de Sitter space every point $p$ of a null geodesic $\gamma$ lies on the bifurcation surface of a Killing horizon generated in part by $\gamma$ (the horizon is the past light cone of the endpoint of $\gamma$ at future infinity, cut by the future light cone of a point of past infinity chosen so that it meets $\gamma$ at $p$), and on that horizon the Killing boost acts on $\gamma$ as the affine dilation about $p$, at rapidity period $2\pi$ in its own normalisation (as in Theorem \claimref{synthesis.positive-null-energy-gives-r-premises}, case (b)). If the state of a free field has geometric modular action for the de Sitter wedges (the modular group of each wedge algebra is its Killing boost at rapidity period $2\pi$, the form in which the Bisognano-Wichmann property holds for the de Sitter-invariant vacua \cite{BorchersBuchholz1999}, cited from memory) and if each wedge algebra is generated by the modes on the future half of its horizon, then on $\gamma$ the hypothesis of Step 4 holds for every $p$, so the state of the modes on $\gamma$ is $\omega_{0}$: (N$_{+}$) case (b) for that field, massive fields included, which the store's conformal-energy route does not reach (Theorem \claimref{synthesis.de-sitter-horizon-translation-is-conformal-null-rotation}). The two hypotheses (geometric modular action, generation by the horizon modes) are not proved here. The period $2\pi$ is part of the first (it is essentially the Bisognano-Wichmann property, from which the null-vacuum restriction on the horizon is already known for de Sitter-invariant vacua), so the derivation of the period in Steps 3 and 4 does no work here; the gain over the stored conformal route is the massive case only, and it rests on the second hypothesis, since for massive fields the restriction to the horizon is not automatic. On the eternal hole only the bifurcation surface carries a Killing boost, so Step 4 does not apply there.

\emph{Scope.} Steps 1 to 5 concern the fermionic (CAR) continuum mode space of one generator; the bosonic analogue needs the uniqueness of CCR KMS states and is not claimed. The hypothesis is natural where every point of a generator is the edge of a Killing wedge: Minkowski null planes (Rindler wedges) and, argued, de Sitter horizons; not the eternal hole. The geometric action of the wedge reflections, which gives spin-statistics for interacting fields, is the companion statement for the modular conjugations (\claimref{synthesis.wedge-reflections-compose-to-the-two-pi-rotor-so-geometric-modular-conjugations-give-interacting-spin-statistics}).

\emph{Status and novelty.} Steps 1 to 5 proved (Step 1 by the stored uniqueness of the I-state, the residue series standard and checked); Step 6 argued, given two unproved hypotheses and a literature result cited from memory. Standard: the I-state kernel; the forward direction (the null vacuum is KMS at $2\pi$ about every point, Theorem \claimref{synthesis.positive-null-energy-gives-r-premises}(ii) with translation invariance); the use of half-sided modular structure to obtain translations is in the spirit of Borchers and Wiesbrock (Fact \claimref{angle.imported-theorems-deliver-their-hypotheses}). New: the converse for arbitrary states, with the period derived rather than assumed, and its reading as a temperature-free equivalent of (N$_{+}$) on one generator; the de Sitter reading. Transport-absent.
